% ============================================================================
% Speichermanagement (IBM): Anleitung fuer den Vorstand
% SD-Karte fuer ein neues Mitglied erstellen, testen und uebergeben.
% Bauen: pdflatex ibm-vorstand-anleitung.tex  (zweimal fuer das Inhaltsverzeichnis)
% Technischer Hintergrund: Batteriemanagement/openhab/setup/README.md
% ============================================================================
\documentclass[a4paper,11pt]{article}
\usepackage[utf8]{inputenc}
\usepackage[T1]{fontenc}
\usepackage[ngerman]{babel}
\usepackage[margin=2.5cm]{geometry}
\usepackage{lmodern}
\usepackage{parskip}
\usepackage{enumitem}
\usepackage{tabularx}
\usepackage{booktabs}
\usepackage{xcolor}
\usepackage[colorlinks=true,linkcolor=black,urlcolor=blue!50!black]{hyperref}

\newcommand{\dash}{Dashboard \url{https://ischlstrom.org/board/openhab}}
\newcommand{\phase}[1]{\textsf{\small #1}}

\begin{document}

\begin{center}
    {\LARGE\bfseries Speichermanagement (IBM)}\\[0.3em]
    {\large Anleitung für den Vorstand: SD-Karte für ein neues Mitglied}\\[0.3em]
    {\small ISCHLSTROM Energiegemeinschaft, Stand: August 2026}
\end{center}

\section*{Worum es geht}

Mitglieder mit Batteriespeicher bekommen von uns einen Raspberry Pi, der
den Speicher im Sinn der Energiegemeinschaft steuert (laden, wenn Strom
in der Gemeinschaft übrig ist). Die gesamte Einrichtung ist automatisiert:
Der Vorstand erstellt am Dashboard eine SD-Karte, testet den Pi kurz im
eigenen Netzwerk und übergibt ihn dann dem Mitglied. Beim Mitglied muss
niemand etwas konfigurieren; der Pi meldet jeden Schritt an ischlstrom.org,
den Fortschritt sehen Vorstand und Mitglied im Browser.

\subsection*{Was man braucht}

\begin{itemize}[nosep]
    \item Vorstandszugang auf \dash
    \item Raspberry Pi 4B mit Netzteil, SD-Karte ab 16 GB
    \item einen Rechner mit SD-Kartenleser und dem Programm
          \emph{Raspberry Pi Imager} (\url{https://www.raspberrypi.com/software/})
          oder \emph{balenaEtcher}
    \item ein freies LAN-Kabel am eigenen Router (für den Test)
\end{itemize}

\section{SD-Karte am Dashboard vorbereiten}

Auf dem \dash{} im Abschnitt \emph{Einrichtung} das Mitglied per
Namenssuche auswählen und \emph{SD-Karte vorbereiten} klicken. Dabei
erzeugt der Server alles Nötige selbst: Passwörter, Fernwartungszugang,
das openHAB-Cloud-Konto \texttt{<nr>@ischlstrom.org} und einen
Provisionierungs-Code (60 Tage gültig).

\textbf{Wichtig:} Dabei gleich das \emph{Wechselrichter-Profil} setzen
(Fronius GEN24, Fronius SnapINverter, Sigenergy, Deye oder Victron; steht
in der Anmeldung des Mitglieds, im Zweifel nachfragen). Ohne Profil
versucht der Pi, den Wechselrichter im Netzwerk zu erkennen; in Ihrem
Testnetz gibt es ihn aber nicht, und die Einrichtung bleibt bei
\phase{Wechselrichter nicht eindeutig} stehen. Das Profil lässt sich
notfalls auch später am Dashboard setzen, die Einrichtung wartet darauf.

Hat das Mitglied kein LAN-Kabel am Aufstellort, können hier auch
WLAN-Name und -Passwort des Mitglieds hinterlegt werden (LAN bleibt die
Empfehlung).

Die Anlage erscheint danach als Karte oben am Dashboard. Alle weiteren
Schritte und Knöpfe (Image, Passwörter, Fortschritt, Profil setzen,
Neuer Code) stehen in der \textbf{Detailansicht der Anlage}, die sich
per Klick auf die Karte öffnet.

\section{Image erstellen und auf die Karte schreiben}

\begin{enumerate}[nosep]
    \item In der Detailansicht der Anlage \emph{Image erstellen} klicken. Der
          Server baut das fertige Karten-Image (aktuelles openHABian plus
          unsere Konfiguration); das dauert einige Minuten. Danach
          \emph{Image herunterladen} (Datei \texttt{pi-<nr>.img.gz},
          etwa 1{,}5 GB).
    \item Raspberry Pi Imager starten: bei \emph{Betriebssystem} ganz
          unten \emph{Eigenes Image verwenden} wählen und die
          heruntergeladene Datei auswählen (entpacken ist nicht nötig).
          SD-Karte wählen und schreiben.
    \item \textbf{Wichtig:} Die Frage \emph{"Möchten Sie OS-Anpassungen
          anwenden?"} mit \textbf{Nein} beantworten. Eigene Anpassungen
          (Hostname, WLAN, Benutzer) würden die mitgelieferte
          Konfiguration überschreiben und die automatische Einrichtung
          verhindern. Mit balenaEtcher stellt sich die Frage nicht.
\end{enumerate}

\section{Testlauf im eigenen Netzwerk}

Die Karte in den Pi stecken, den Pi per LAN-Kabel mit dem eigenen Router
verbinden und das Netzteil anstecken. Mehr ist nicht zu tun; nicht
zwischendurch den Strom trennen. Der Pi installiert zuerst das
Betriebssystem und openHAB (der längste Teil), dann richtet er sich bei
ischlstrom.org ein. Insgesamt dauert das 30 bis 45 Minuten.

Den Fortschritt zeigen die Karte am Dashboard und die Detailansicht der
Anlage an (die Seiten aktualisieren sich selbst), ebenso die Mitgliedsseite
\emph{Speichermanagement}. Die Phasen laufen von
\phase{Konfiguration geladen} über \phase{Fernwartung},
\phase{Passwörter}, \phase{Cloud-Verbindung} bis \phase{Steuerung} und
\phase{Oberfläche}.

\subsection*{Erwartetes Ergebnis des Tests}

Da es Ihren Testaufbau ohne Wechselrichter gibt, endet die Einrichtung
absichtlich \textbf{nicht} auf "fertig", sondern bei

\begin{center}
    \phase{Wartet, wird automatisch fortgesetzt}
\end{center}

(gelber Balken). Das ist der gewollte Endzustand: Alles außer dem
Wechselrichter ist eingerichtet und geprüft; den Wechselrichter findet
der Pi erst beim Mitglied und macht dann von selbst weiter. Ein roter
Balken (\phase{Fehler bei: \ldots}) ist dagegen ein echtes Problem,
siehe Fehlerbehebung.

Zwei Dinge kurz prüfen:

\begin{itemize}[nosep]
    \item Die Anlage zeigt am Dashboard einen Fernwartungs-Haken
          (WireGuard verbunden) und "zuletzt gemeldet" mit frischer
          Uhrzeit.
    \item Anmeldung auf \url{https://remote.hac.ischlstrom.org} mit dem
          Cloud-Konto der Anlage (Benutzer und Passwort stehen in der
          Detailansicht der Anlage unter \emph{Passwörter anzeigen}): die
          openHAB-Oberfläche des Pi erscheint.
\end{itemize}

Danach den Pi vom Strom trennen (kurz warten, bis die grüne LED ruhig
ist, oder vorher über die Cloud-Oberfläche herunterfahren).

\section{Übergabe an das Mitglied}

Der Pi kommt beim Mitglied \textbf{in dasselbe Netzwerk wie der
Wechselrichter} (LAN-Kabel am Router bzw. das hinterlegte WLAN) und ans
Netzteil. Der Pi setzt die Einrichtung selbst fort: Er sucht alle
10 Minuten nach dem Wechselrichter, richtet ihn ein und meldet am Ende
\phase{Einrichtung abgeschlossen}.

\textbf{Nur beim Fronius GEN24}: Das Mitglied muss sein
Wechselrichter-Passwort (Benutzer \texttt{customer} der
Weboberfläche des GEN24) eintragen, und zwar auf seiner Seite
\emph{Speichermanagement} auf ischlstrom.org (anmelden, Anlage
auswählen, Feld \emph{Passwort des Wechselrichters}). Das kann vor oder
nach dem Anstecken passieren; der Pi holt es sich automatisch. Das
Passwort wird verschlüsselt übertragen, einmalig an den Pi ausgeliefert
und danach am Server gelöscht. Die anderen Profile brauchen kein
Passwort.

Dem Mitglied zeigen bzw. mitgeben:

\begin{itemize}[nosep]
    \item Die Seite \emph{Speichermanagement} im Mitgliederbereich:
          Fortschritt, Batterie-Ladestand, Cloud-Zugangsdaten.
    \item Zugriff per openHAB-App (App Store / Play Store): als
          Remote-URL \texttt{https://hac.ischlstrom.org} eintragen und
          mit dem Cloud-Konto anmelden; im Browser
          \url{https://remote.hac.ischlstrom.org}.
    \item Die \textbf{Einwilligung} zur Datennutzung: Das Mitglied stimmt
          auf seiner Seite \emph{Speichermanagement} online zu (das
          Dashboard zeigt je Anlage ein Badge \emph{Einwilligung}); ohne
          Einwilligung sieht das Mitglied dort nur den Text und keine
          Anlagendaten. Wer nicht online zustimmen mag: Papier-Variante
          \texttt{speichermanagement-einverstaendnis.pdf} in diesem
          Ordner unterschreiben lassen und ablegen.
\end{itemize}

\section{Fehlerbehebung}

\begin{tabularx}{\textwidth}{@{}p{0.34\textwidth}X@{}}
    \toprule
    \textbf{Anzeige am Dashboard} & \textbf{Was tun} \\
    \midrule
    \phase{Wechselrichter nicht eindeutig} &
    Wechselrichter-Profil am Dashboard setzen (\emph{Profil setzen});
    die Einrichtung macht von selbst weiter. \\
    \addlinespace
    \phase{Wartet auf Wechselrichter-Passwort} &
    Nur GEN24: Mitglied (oder Vorstand am Dashboard) trägt das Passwort
    ein; der Pi fragt alle zwei Minuten nach. \\
    \addlinespace
    \phase{Fehler bei: \ldots} (roter Balken) &
    Einrichtung ist an diesem Schritt gescheitert. Kurz warten (der Pi
    wiederholt vieles selbst); bleibt es rot, an den technischen
    Betrieb melden (info@ischlstrom.org). \\
    \addlinespace
    Keine Phase, obwohl der Pi seit über einer Stunde läuft &
    Netzwerkverbindung des Pi prüfen (LAN-Kabel, Router-Geräteliste).
    Hinweis "mit einem alten oder abgelaufenen Code gebaut" am
    Dashboard: \emph{Neuer Code}, dann \emph{Image erstellen} und die
    Karte neu schreiben. \\
    \addlinespace
    Cloud-Konto oder Mail-Alias mit Fehler &
    Knopf \emph{Cloud-Konto erneut} bzw. \emph{Mail-Alias erneut}
    klicken. \\
    \addlinespace
    SD-Karte defekt &
    \emph{Neuer Code}, \emph{Image erstellen}, neue Karte schreiben;
    alle Einstellungen liegen am Server, es geht nichts verloren. Beim
    GEN24 muss das Mitglied sein Wechselrichter-Passwort neu eintragen. \\
    \bottomrule
\end{tabularx}

\vspace{1em}
\begin{center}
    \small Technische Details: \texttt{Batteriemanagement/openhab/setup/README.md}
    im Repository \,\textbullet\, Fragen: info@ischlstrom.org
\end{center}

\end{document}
